\chapter{Exclusión Mutua y Mecanismos de Sincronización Básicos}
\label{cap:tema2_parte1}

\cuadrofuentes{
    \textbf{Fuentes utilizadas:} Transparencias Oficiales del Profesor (Manuel I. Capel Tuñón), Temario Completo de Infomates (LosDelDGIIM) y Apuntes de Ismael Sallami Moreno.
}

% ==============================================================================
% 2.1. EL PROBLEMA DE LA EXCLUSIÓN MUTUA
% ==============================================================================
\section{El Problema de la Exclusión Mutua}

Cuando dos o más procesos o hebras se ejecutan concurrentemente y comparten datos o recursos en memoria, surge la necesidad de coordinar sus accesos para evitar que interfieran mutuamente de manera destructiva.

\begin{definicion}[Sección Crítica (SC) y Sección No Crítica (SNC)]
\begin{itemize}[leftmargin=*]
    \item \textbf{Sección Crítica (SC):} Es la porción de código de un proceso en la cual se accede y manipula uno o más recursos o variables compartidas. Para evitar inconsistencias de datos, debe cumplirse la propiedad de que en cualquier instante de tiempo a lo sumo un solo proceso se encuentre ejecutando su sección crítica.
    \item \textbf{Sección No Crítica (SNC):} Cualquier porción del código en la que el proceso trabaja únicamente con variables locales privadas o recursos no compartidos. Las acciones en la SNC no interfieren con otros procesos.
\end{itemize}
\end{definicion}

\subsection{Estructura Canónica de un Proceso Concurrente}
Cualquier proceso que compita por una sección crítica responde al siguiente esquema cíclico:

\begin{lstlisting}[language=C++]
void Proceso(int i) {
    while (true) {
        // 1. Protocolo de entrada: solicita permiso y espera si la SC esta ocupada
        protocolo_entrada(i);
        
        // 2. Seccion Critica (SC): acceso seguro a variables compartidas
        seccion_critica();
        
        // 3. Protocolo de salida: libera el recurso y avisa a otros procesos
        protocolo_salida(i);
        
        // 4. Seccion No Critica (SNC): computo independiente
        seccion_no_critica();
    }
}
\end{lstlisting}

\subsection{Las Cuatro Condiciones de Dijkstra}
Edsger W. Dijkstra formuló en 1965 los cuatro requisitos formales indispensables que cualquier protocolo de entrada y salida debe satisfacer obligatoriamente para considerarse una solución válida y correcta al problema de la exclusión mutua:

\begin{definicion}[Condiciones de Dijkstra]
\begin{enumerate}[leftmargin=*]
    \item \textbf{Exclusión Mutua Estricta (Seguridad):} Nunca puede haber más de un proceso simultáneamente dentro de la sección crítica ($N_{SC} \le 1$).
    \item \textbf{No Suposiciones sobre el Hardware ni Velocidades:} No se puede asumir ninguna hipótesis sobre el número de procesadores, ni sobre la velocidad relativa de ejecución de los procesos.
    \item \textbf{Progreso / Ausencia de Interbloqueo (Seguridad/Vivacidad):} Un proceso que se encuentre ejecutando en su sección no crítica (SNC) o que haya terminado no puede impedir que otros procesos interesados accedan a la sección crítica. Si ningún proceso está en la SC y hay procesos intentando entrar, la decisión de cuál entra debe tomarse en un tiempo finito.
    \item \textbf{Espera Limitada / Ausencia de Inanición (Vivacidad):} Ningún proceso que desee entrar a la sección crítica debe ser postergado indefinidamente. El tiempo de espera para entrar debe ser finito.
\end{enumerate}
\end{definicion}

% ==============================================================================
% 2.2. ALGORITMOS SOFTWARE CON ESPERA ACTIVA
% ==============================================================================
\section{Algoritmos Software con Espera Activa (Spinlocks)}

En los inicios de la computación, antes de que los sistemas operativos y los procesadores ofrecieran instrucciones de sincronización dedicadas, se investigó cómo resolver la exclusión mutua exclusivamente mediante instrucciones de lectura y escritura ordinarias sobre memoria compartida.

\subsection{El Método del Refinamiento Sucesivo: Los Cuatro Intentos}
Para comprender la lógica de los algoritmos correctos, es clásico analizar los cuatro intentos fallidos para 2 procesos ($P_0$ y $P_1$):

\subsubsection{Intento 1: Alternancia Estricta (Variable Turno)}
\begin{lstlisting}[language=C++]
int turno = 0; // Variable compartida

// Proceso P0:
while (turno != 0) { /* Espera activa */ }
seccion_critica();
turno = 1;

// Proceso P1:
while (turno != 1) { /* Espera activa */ }
seccion_critica();
turno = 0;
\end{lstlisting}
\begin{itemize}
    \item \textbf{Cumple:} Exclusión mutua estricta.
    \item \textbf{Falla:} \textbf{Viola la condición de progreso}. Si $P_0$ ejecuta su SC, pone \texttt{turno = 1}, pero $P_1$ se queda permanentemente en su sección no crítica, $P_0$ jamás podrá volver a entrar a la SC porque \texttt{turno} sigue siendo 1. La velocidad de un proceso queda atada a la del otro.
\end{itemize}

\subsubsection{Intento 2: Banderas de Estado}
Se sustituye el turno por dos variables booleanas de presencia: \texttt{bool en\_sc[2] = \{false, false\};}
\begin{lstlisting}[language=C++]
// Proceso Pi:
while (en_sc[1 - i]) { /* Espera a que el otro salga */ }
en_sc[i] = true;
seccion_critica();
en_sc[i] = false;
\end{lstlisting}
\begin{itemize}
    \item \textbf{Falla:} \textbf{Viola la exclusión mutua}. Si ambos procesos evalúan el bucle \texttt{while} a la vez cuando ambas banderas son \texttt{false}, ambos salen del bucle, ambos escriben \texttt{true} y ambos entran a la SC a la vez.
\end{itemize}

\subsubsection{Intento 3: Banderas de Intención Previas}
Para evitar el fallo del Intento 2, cada proceso declara su intención de entrar \textit{antes} de comprobar al otro:
\begin{lstlisting}[language=C++]
bool quiere_entrar[2] = {false, false};

// Proceso Pi:
quiere_entrar[i] = true;
while (quiere_entrar[1 - i]) { /* Espera */ }
seccion_critica();
quiere_entrar[i] = false;
\end{lstlisting}
\begin{itemize}
    \item \textbf{Falla:} \textbf{Sufre de Interbloqueo (Deadlock)}. Si $P_0$ y $P_1$ ejecutan \texttt{quiere\_entrar[i] = true} simultáneamente, ambos encuentran la bandera del vecino en \texttt{true} y ambos entran en bucle infinito de espera activa, bloqueándose para siempre.
\end{itemize}

\subsubsection{Intento 4: Retirada de Cortesía}
Para solucionar el interbloqueo del Intento 3, si un proceso ve que el otro también quiere entrar, baja su bandera momentáneamente para cederle el paso:
\begin{lstlisting}[language=C++]
// Proceso Pi:
quiere_entrar[i] = true;
while (quiere_entrar[1 - i]) {
    quiere_entrar[i] = false;
    sleep(random);
    quiere_entrar[i] = true;
}
seccion_critica();
quiere_entrar[i] = false;
\end{lstlisting}
\begin{itemize}
    \item \textbf{Falla:} \textbf{Sufre de Bloqueo Activo (Livelock)}. Si ambos procesos ejecutan a la par, pueden levantar y bajar sus banderas al unísono infinitamente sin que ninguno llegue a entrar jamás a la SC.
\end{itemize}

% ------------------------------------------------------------------------------
\subsection{El Algoritmo de Dekker (1965)}
Fue el primer algoritmo completamente correcto conocido en la historia de la computación para resolver el problema de la exclusión mutua para 2 procesos. Combina las banderas de intención del Intento 3 con una variable de desempate (\texttt{turno}) para romper la simetría cuando ambos desean entrar a la vez.

\begin{algoritmo}[Algoritmo de Dekker para 2 Procesos]
\begin{lstlisting}[language=C++]
bool quiere[2] = {false, false};
int turno = 0;

void Proceso(int i) {
    int j = 1 - i; // Identificador del otro proceso
    while (true) {
        quiere[i] = true;
        while (quiere[j]) {
            if (turno != i) {
                quiere[i] = false;
                while (turno != i) { /* Espera pasiva de su turno */ }
                quiere[i] = true;
            }
        }
        
        seccion_critica();
        
        turno = j;          // Cede el turno al otro proceso
        quiere[i] = false;  // Indica que ya ha salido
        
        seccion_no_critica();
    }
}
\end{lstlisting}
\end{algoritmo}

% ------------------------------------------------------------------------------
\subsection{El Algoritmo de Peterson para 2 Procesos (1981)}
Gary L. Peterson simplificó drásticamente la solución de Dekker en 1981 con un algoritmo sumamente elegante, considerado una obra maestra de la programación concurrente.

\begin{algoritmo}[Algoritmo de Peterson para 2 Procesos]
\begin{lstlisting}[language=C++]
bool quiere[2] = {false, false};
int turno = 0;

void Proceso(int i) {
    int j = 1 - i;
    while (true) {
        // PROTOCOLO DE ENTRADA:
        quiere[i] = true;  // 1. Declaro mi intencion de entrar
        turno = j;         // 2. Cortesamente ofrezco el turno al otro
        while (quiere[j] && turno == j) {
            // Espera activa mientras el otro quiera Y mantenga el turno
        }
        
        seccion_critica();
        
        // PROTOCOLO DE SALIDA:
        quiere[i] = false; // Indico que he salido de la SC
        
        seccion_no_critica();
    }
}
\end{lstlisting}
\end{algoritmo}

\subsubsection{Demostración Formal de la Exclusión Mutua en Peterson}
Demostremos que $P_0$ y $P_1$ nunca pueden estar juntos en la SC por reducción al absurdo:
\begin{proof}
Supongamos que ambos procesos están en la SC simultáneamente. Para que $P_0$ haya salido de su bucle de espera, debe haber evaluado como falsa la condición:
\[
\text{quiere}[1] == \text{false} \quad \lor \quad \text{turno} == 0
\]
Como ambos están en la SC, obligatoriamente $\text{quiere}[0] == \text{true}$ y $\text{quiere}[1] == \text{true}$. Por consiguiente, la única forma de que $P_0$ entrara es que viera $\text{turno} == 0$.
Análogamente, para que $P_1$ haya entrado, debe haber visto $\text{turno} == 1$.
Sin embargo, \texttt{turno} es una variable atómica escalar que solo puede tener un único valor (0 o 1). El valor de \texttt{turno} corresponde a la **última asignación realizada**.
Si $P_0$ asignó $\text{turno} = 1$ en último lugar, entonces $\text{turno} == 1$, lo que impidió entrar a $P_0$ y permitió entrar solo a $P_1$. Es imposible que $\text{turno} == 0$ y $\text{turno} == 1$ simultáneamente. Contradicción. Por tanto, se garantiza la exclusión mutua.
\end{proof}

% ------------------------------------------------------------------------------
\subsection{Exclusión Mutua para N Procesos}
Para extender la exclusión mutua por software a $N$ procesos existen dos algoritmos célebres:

\begin{itemize}[leftmargin=*]
    \item \textbf{Filtro de Peterson:} Generaliza el algoritmo para 2 procesos creando $N-1$ etapas o niveles de competición. En cada nivel, al menos un proceso queda esperando en el filtro, de forma que a la etapa $N-1$ solo llega un único ganador.
    \item \textbf{Algoritmo de la Panadería de Lamport (\textit{Bakery Algorithm}):} Simula el sistema de turnos de una panadería:
    \begin{enumerate}
        \item Cada proceso $i$ que desea entrar toma un número de turno: $\text{turno}[i] = 1 + \max(\text{turno}[0], \dots, \text{turno}[n-1])$.
        \item Se atiende a los procesos por orden lexicográfico estricto del par $(\text{turno}[i], i)$.
        \item \textbf{Propiedad excepcional:} Es correcto incluso si las lecturas y escrituras en memoria no son estrictamente atómicas (tolera lecturas solapadas con escrituras en variables compartidas).
    \end{enumerate}
\end{itemize}

% ==============================================================================
% 2.3. SOPORTE HARDWARE PARA EXCLUSIÓN MUTUA
% ==============================================================================
\section{Soporte Hardware para Exclusión Mutua}

Aunque los algoritmos de software puros son teóricamente fascinantes, en la práctica presentan graves inconvenientes: consumen el 100\% de CPU en bucles de espera activa (\textit{busy waiting}), son complejos para $N$ procesos y en los procesadores modernos con ejecución fuera de orden (\textit{out-of-order execution}) requieren barreras de memoria explícitas (\textit{memory fences}).

Por ello, los procesadores modernos incorporan instrucciones atómicas a nivel de instrucción máquina:

\subsection{1. Deshabilitación de Interrupciones}
En sistemas monoprocesador clásicos, un proceso de nivel núcleo puede deshabilitar las interrupciones hardware antes de entrar a la SC y rehabilitarlas al salir. Sin embargo:
\begin{itemize}
    \item Es inaceptable en espacio de usuario (un proceso con bucle infinito congelaría el sistema entero).
    \item Es completamente inútil en sistemas multiprocesador multinúcleo (deshabilitar interrupciones en el Core 1 no impide que el Core 2 acceda a la misma memoria RAM).
\end{itemize}

\subsection{2. Instrucciones Atómicas Read-Modify-Write (RMW)}
Estas instrucciones ejecutan una lectura y una modificación de una posición de memoria de forma indivisible a nivel de bus de memoria o caché:

\begin{definicion}[Test-and-Set (TAS)]
Lee el valor original de una variable booleana en memoria y escribe incondicionalmente \texttt{true}, devolviendo el valor anterior de forma atómica:
\begin{lstlisting}[language=C++]
bool TestAndSet(bool &cerrojo) {
    bool anterior = cerrojo;
    cerrojo = true;
    return anterior;
}
\end{lstlisting}
Implementación de un cerrojo simple de exclusión mutua (\textit{Spinlock}):
\begin{lstlisting}[language=C++]
bool cerrojo = false; // Libre

void entrar() {
    while (TestAndSet(cerrojo)) { /* Espera activa */ }
}
void salir() {
    cerrojo = false; // Liberacion atomica
}
\end{lstlisting}
\end{definicion}

\begin{definicion}[Compare-and-Swap (CAS)]
Instrucción más potente y flexible. Compara el contenido de una dirección de memoria con un valor esperado; solo si coinciden, almacena el nuevo valor:
\begin{lstlisting}[language=C++]
bool CompareAndSwap(int *dir, int valor_esperado, int nuevo_valor) {
    if (*dir == valor_esperado) {
        *dir = nuevo_valor;
        return true;
    }
    return false;
}
\end{lstlisting}
El CAS es la piedra angular de las estructuras de datos concurrentes libres de bloqueos (\textit{lock-free programming}) y de la biblioteca \texttt{<atomic>} de C++11 (como \texttt{std::atomic<T>::compare\_exchange\_weak}).
\end{definicion}

% ==============================================================================
% 2.4. SEMÁFOROS DE DIJKSTRA
% ==============================================================================
\section{Semáforos de Dijkstra}

Para superar el grave problema del desperdicio de CPU que produce la espera activa, Edsger Dijkstra introdujo en 1965 el concepto de **Semáforo**.

\begin{definicion}[Semáforo]
Un semáforo $S$ es una variable protegida especial con valor entero que dispone de una cola de espera de procesos bloqueados y sobre la que **únicamente se permiten tres operaciones**:
\begin{enumerate}
    \item \textbf{Inicialización:} Establece el valor inicial no negativo de $S$.
    \item \textbf{Operación P(S) / wait(S):} (Del holandés \textit{proberen}, probar/decrementar). Decrementa el valor del semáforo. Si el valor resultante es negativo (o si valía 0), el proceso invocador suspende su ejecución y pasa a la cola de bloqueados del semáforo.
    \item \textbf{Operación V(S) / signal(S) / post(S):} (Del holandés \textit{verhogen}, incrementar). Incrementa el valor del semáforo. Si hay procesos esperando en la cola, despierta a uno de ellos y lo pasa a estado listo.
\end{enumerate}
Tanto $P(S)$ como $V(S)$ son operaciones **estrictamente atómicas**.
\end{definicion}

\subsection{Tipos de Semáforos}
\begin{itemize}[leftmargin=*]
    \item \textbf{Semáforo Binario (Mutex):} Su valor solo puede ser 0 o 1. Se utiliza típicamente para exclusión mutua (inicializado a 1).
    \item \textbf{Semáforo General o Contador:} Su valor es un entero mayor o igual que 0. Se inicializa al número $K$ de unidades disponibles de un recurso compartido.
\end{itemize}

\subsection{El Invariante del Semáforo}
Sea $S_{init}$ el valor inicial del semáforo, $N_P$ el número de operaciones $P$ completadas con éxito y $N_V$ el número de operaciones $V$ completadas. Se verifica siempre la siguiente propiedad invariante fundamental:
\begin{equation}
S_{actual} = S_{init} + N_V - N_P \ge 0
\end{equation}
Por tanto, si un semáforo se inicializa a 1 (para exclusión mutua):
\begin{equation}
N_P - N_V \le 1
\end{equation}
Como los procesos en la SC corresponden a los que han completado $P$ y no han hecho $V$, el número de procesos en la SC cumple $N_{SC} = N_P - N_V \le 1$, garantizando formalmente la exclusión mutua.

\subsection{Semáforos en la Práctica: POSIX y C++20}
En entornos UNIX/Linux se utiliza la API POSIX (\texttt{semaphore.h}):
\begin{lstlisting}[language=C++,caption={Uso de semáforos POSIX en C++}]
#include <semaphore.h>

sem_t sem_mutex;

int main() {
    sem_init(&sem_mutex, 0, 1); // Inicializa a 1 (compartido entre hebras)
    
    // Protocolo de entrada:
    sem_wait(&sem_mutex); // Equivale a P()
    
    // Seccion Critica...
    
    // Protocolo de salida:
    sem_post(&sem_mutex); // Equivale a V()
    
    sem_destroy(&sem_mutex);
}
\end{lstlisting}

\subsection{Limitaciones de los Semáforos}
Aunque los semáforos son muy potentes y resuelven el problema de la espera activa, su uso en proyectos grandes resulta muy propenso a errores:
\begin{itemize}
    \item \textbf{Falta de Estructura:} Son variables globales dispersas por el código; no existe asociación sintáctica entre el semáforo y los datos que protege.
    \item \textbf{Fragilidad:} Un error de programación en un solo proceso (olvidar un \texttt{sem\_post}, invertir el orden de dos \texttt{sem\_wait} o duplicar una llamada) provoca inmediatamente interbloqueo o violación de la exclusión mutua en todo el sistema.
    \item \textbf{Doble Propósito Confuso:} Se usan tanto para exclusión mutua como para sincronización condicional, haciendo el código difícil de mantener y verificar.
\end{itemize}

% ==============================================================================
% 2.5. FUNDAMENTOS DE LOS MONITORES
% ==============================================================================
\section{Introducción a los Monitores como Mecanismo de Alto Nivel}

Para superar las deficiencias de los semáforos, C.A.R. Hoare (1974) y Per Brinch Hansen diseñaron los **Monitores**: un mecanismo de sincronización estructurado de alto nivel integrado en el propio lenguaje de programación.

\begin{definicion}[Monitor]
Un monitor es un Tipo de Dato Abstracto (TDA) concurrente que encapsula:
\begin{enumerate}
    \item Las variables compartidas que definen el estado del recurso protegido (variables permanentes privadas).
    \item Un conjunto de procedimientos públicos accesibles por los procesos externos.
    \item Un bloque de inicialización de las variables internas.
\end{enumerate}
\textbf{Regla de Oro del Monitor:} El propio compilador o entorno de ejecución garantiza de forma automática e implícita la **exclusión mutua** en el acceso a sus procedimientos. En ningún momento puede haber más de un proceso ejecutándose activamente dentro del monitor.
\end{definicion}

\subsection{Estructura Interna de un Monitor}
Todo monitor dispone de dos niveles de colas:
\begin{itemize}
    \item \textbf{Cola de Entrada al Monitor:} Donde esperan los procesos que intentan invocar un procedimiento del monitor mientras otro proceso ya está dentro de él.
    \item \textbf{Variables Condición (\texttt{cond}):} Colas internas donde se suspenden los procesos que ya están dentro del monitor pero que no pueden continuar porque no se cumple una condición lógica sobre las variables del monitor.
\end{itemize}

\subsection{Variables Condición y sus Primitivas}
Una variable de tipo \texttt{cond} (o \texttt{condition\_variable}) no guarda un valor booleano ni numérico; representa exclusivamente una **cola de espera**.
Sobre una variable condición $c$ se definen dos operaciones básicas:
\begin{itemize}
    \item \texttt{c.wait():} El proceso que la invoca se bloquea en la cola de $c$ y **libera automáticamente la exclusión mutua del monitor**, permitiendo que otro proceso entre a modificar el estado.
    \item \texttt{c.signal():} Despierta a uno de los procesos suspendidos en la cola de $c$. Si no hay ningún proceso esperando en $c$, la llamada a \texttt{signal()} no tiene efecto (se pierde, a diferencia de un semáforo que incrementaría su valor).
\end{itemize}

% ==============================================================================
% 2.6. SEMÁNTICAS DE SEÑALIZACIÓN EN MONITORES
% ==============================================================================
\section{Semánticas de Señalización en Monitores}

El dilema fundamental en el diseño de un monitor ocurre en el instante exacto en que un proceso $P$ dentro del monitor ejecuta \texttt{c.signal()} y despierta al proceso suspendido $Q$. En ese instante habría **dos procesos dentro del monitor** ($P$ y $Q$), lo que violaría la exclusión mutua implícita.

Para resolver este conflicto, existen cuatro semánticas principales:

\begin{table}[h!]
\centering
\small
\renewcommand{\arraystretch}{1.3}
\begin{tabular}{@{}lp{4.2cm}p{4.2cm}c@{}}
\toprule
\textbf{Semántica} & \textbf{Acción sobre el señalador ($P$)} & \textbf{Acción sobre el despertado ($Q$)} & \textbf{¿Desplazante?} \\ \midrule
\textbf{SU} (Señalar y Espera Urgente) & Pasa a una cola prioritaria de procesos urgentes. & Entra inmediatamente al monitor y reanuda su ejecución. & \textbf{Sí} (Hoare) \\
\textbf{SC} (Señalar y Continuar) & Continúa ejecutando dentro del monitor hasta salir o esperar. & Pasa a la cola general de entrada o lista de listos. & \textbf{No} (Mesa/Java) \\
\textbf{SS} (Señalar y Salir) & Debe salir del monitor en ese mismo instante. & Entra inmediatamente al monitor. & \textbf{Sí} \\
\textbf{SE} (Señalar y Esperar) & Pasa a la cola ordinaria de entrada. & Reanuda la ejecución dentro del monitor. & \textbf{Sí} \\ \bottomrule
\end{tabular}
\caption{Comparativa de semánticas de señalización en monitores.}
\end{table}

\begin{aviso}[La Regla de Oro: \texttt{if} vs. \texttt{while} al evaluar la condición]
\begin{itemize}[leftmargin=*]
    \item En la semántica **SU (Hoare)**, el proceso despertado se ejecuta de inmediato sin que nadie modifique el estado entre la señal y su reanudación. Por tanto, basta con comprobar la condición con un condicional:
    \begin{lstlisting}[language=C++]
    if (!condicion_cumplida) c.wait(); // Valido en semantica SU (Hoare)
    \end{lstlisting}
    \item En la semántica **SC (Mesa / Hansen / C++ / Java)**, otros procesos pueden intercalarse y consumir el recurso antes de que el proceso despertado consiga retomar el monitor. Por ello, **es estrictamente obligatorio reevaluar la condición en un bucle \texttt{while}**:
    \begin{lstlisting}[language=C++]
    while (!condicion_cumplida) c.wait(); // OBLIGATORIO en semantica SC
    \end{lstlisting}
\end{itemize}
\end{aviso}

% ==============================================================================
% 2.7. IMPLEMENTACIÓN DE MONITORES CON SEMÁFOROS
% ==============================================================================
\section{Implementación de Monitores usando Semáforos}

Para entender cómo el sistema operativo o el runtime de un lenguaje materializa un monitor (con semántica SU de Hoare), se utiliza la construcción clásica de C.A.R. Hoare basada en semáforos:

\begin{algoritmo}[Esquema de Hoare para Implementar Monitores con Semáforos]
Se definen las siguientes variables de sincronización internas:
\begin{lstlisting}[language=C++]
sem_t mutex;       // Controla la entrada al monitor (inicializado a 1)
sem_t urgente;     // Cola para procesos que cedieron el monitor por signal (inic. a 0)
int num_urgentes = 0; // Contador de procesos en la cola urgente

// Por cada variable condicion c:
sem_t c_sem;       // Cola del semaforo para c.wait() (inic. a 0)
int c_count = 0;   // Contador de procesos esperando en c
\end{lstlisting}

\textbf{1. Entrada y Salida de cada Procedimiento del Monitor:}
\begin{lstlisting}[language=C++]
void Procedimiento_Monitor() {
    // ENTRADA AL PROCEDIMIENTO:
    sem_wait(&mutex);
    
    // ... Codigo del procedimiento ...
    
    // SALIDA DEL PROCEDIMIENTO:
    // Si hay procesos en la cola urgente, tienen prioridad sobre los de entrada:
    if (num_urgentes > 0) {
        sem_post(&urgente);
    } else {
        sem_post(&mutex);
    }
}
\end{lstlisting}

\textbf{2. Implementación de \texttt{c.wait()}:}
\begin{lstlisting}[language=C++]
void wait_c() {
    c_count++;
    // Si habia procesos urgentes, cede a uno urgente; si no, libera la entrada:
    if (num_urgentes > 0) {
        sem_post(&urgente);
    } else {
        sem_post(&mutex);
    }
    sem_wait(&c_sem); // Se bloquea en la condicion
    c_count--;
}
\end{lstlisting}

\textbf{3. Implementación de \texttt{c.signal()} (Semántica SU):}
\begin{lstlisting}[language=C++]
void signal_c() {
    if (c_count > 0) {
        num_urgentes++;
        sem_post(&c_sem);    // Despierta al proceso en la condicion
        sem_wait(&urgente);  // El proceso senalador se bloquea en urgentes
        num_urgentes--;
    }
}
\end{lstlisting}
\end{algoritmo}
Este esquema demuestra de forma transparente cómo se garantiza tanto la exclusión mutua del monitor como la prioridad absoluta de la semántica urgente de Hoare.
